\par\medskip
\noindent\textbf{Corolario 4.13.}
Sea $X_0$ un esquema de tipo finito sobre $k$, liso fuera de un
subconjunto finito
\[
X_0^{\mathrm{Sing}}=\{x_1,\ldots,x_r\},
\]
y supongamos que $H^2(X_0,\underline D^0_{X_0})=0$. Entonces:

\smallskip
\noindent\textbf{(a)} El funtor
\[
\underline M_{X_0}\longrightarrow
\prod_{i=1}^r
\underline M_{\operatorname{Spec}(\widehat{\mathcal O}_{X_0,x_i})}
\]
es liso. En particular,
\[
\dim\underline M_{X_0}
=\dim_k H^1(X_0,\underline D^0_{X_0})
+\sisi{\sum_{i=1}^r
\dim\underline M_{\operatorname{Spec}(\widehat{\mathcal O}_{X_0,x_i})}}{\sum_{c=1}^r
\dim\underline M_{\operatorname{Spec}(\widehat{\mathcal O}_{X_0,x_c})}}.
\]

\smallskip
\noindent\textbf{(b)} Si, además, $X_0$ es localmente una
intersección completa, entonces $\underline M_{X_0}$ es liso.

%%%%%%%%%%  scan idx 120  |  folio 109  %%%%%%%%%%

\smallskip
\noindent\textit{Demostración.}
\textbf{(a)} La primera afirmación se deduce de 4.10 y 4.9(b).
Sean $(R,\mathfrak X)$ y
$(R_1,\mathfrak X_1),\ldots,(R_r,\mathfrak X_r)$ deformaciones
formales versales de $X_0$ y de
\[
\operatorname{Spec}(\mathcal O_{X_0,x_1}),\ldots,
\operatorname{Spec}(\mathcal O_{X_0,x_r}),
\]
respectivamente. Obtenemos un diagrama conmutativo
\[
\begin{tikzcd}[column sep=large,row sep=large]
\underline h_R
  \arrow[r]
  \arrow[d]
& \underline h_{R_1\widehat\otimes_\Lambda\cdots
  \widehat\otimes_\Lambda R_r}
  \arrow[d] \\
\underline M_{X_0}
  \arrow[r]
& \displaystyle\prod_{i=1}^r
  \underline M_{\operatorname{Spec}(\widehat{\mathcal O}_{X_0,x_i})}.
\end{tikzcd}
\]
El funtor compuesto
\[
\underline h_R\longrightarrow\underline M_{X_0}
\longrightarrow
\prod_{i=1}^r
\underline M_{\operatorname{Spec}(\widehat{\mathcal O}_{X_0,x_i})}
\]
es liso, mientras que
\[
\underline h_{R_1\widehat\otimes_\Lambda\cdots
\widehat\otimes_\Lambda R_r}
\longrightarrow
\prod_{i=1}^r
\underline M_{\operatorname{Spec}(\widehat{\mathcal O}_{X_0,x_i})}
\]
es mínimamente liso. Se deduce de 2.7(a) que
\[
\underline h_R\longrightarrow
\underline h_{R_1\widehat\otimes_\Lambda\cdots
\widehat\otimes_\Lambda R_r}
\]
es liso; equivalentemente, $R$ es un anillo de series formales sobre
$R_1\widehat\otimes_\Lambda\cdots\widehat\otimes_\Lambda R_r$.

Por 4.4, la sucesión
\[
0\longrightarrow H^1(X_0,\underline D^0_{X_0})
\longrightarrow[\underline M_{X_0}(k[\varepsilon])]
\longrightarrow
\prod_{i=1}^r
[\underline M_{\operatorname{Spec}(\widehat{\mathcal O}_{X_0,x_i})}
  (k[\varepsilon])]
\]
es exacta. Equivalentemente,
\[
0\longrightarrow H^1(X_0,\underline D^0_{X_0})
\longrightarrow\operatorname{Hom}_{\Lambda\text{-alg}}(R,k[\varepsilon])
\longrightarrow
\operatorname{Hom}_{\Lambda\text{-alg}}
 (R_1\widehat\otimes_\Lambda\cdots\widehat\otimes_\Lambda R_r,
  k[\varepsilon])
\]
es exacta. En consecuencia, $R$ es un anillo de series formales sobre
$R_1\widehat\otimes_\Lambda\cdots\widehat\otimes_\Lambda R_r$ en
$d$ variables, donde
\[
d=\dim_k H^1(X_0,\underline D^0_{X_0}).
\]
Por tanto,
\begin{align*}
\dim\underline M_{X_0}
 &=\dim(k\otimes_\Lambda R) \\
 &=d+\dim\bigl(k\otimes_\Lambda
   (R_1\widehat\otimes_\Lambda\cdots\widehat\otimes_\Lambda R_r)
   \bigr) \\
 &=d+\dim\bigl((k\otimes_\Lambda R_1)
   \widehat\otimes_k\cdots\widehat\otimes_k
   (k\otimes_\Lambda R_r)\bigr) \\
 &=d+\sum_{i=1}^r\dim(k\otimes_\Lambda R_i) \\
 &=d+\sum_{i=1}^r
   \dim\underline M_{\operatorname{Spec}
   (\widehat{\mathcal O}_{X_0,x_i})}.
\end{align*}

%%%%%%%%%%  scan idx 121  |  folio 110  %%%%%%%%%%

\noindent\textbf{(b)} Esto se deduce de 4.10 y 4.9(a).
\hfill$\square$
